\documentclass[12pt,a4paper]{article}

% --- Codificação e língua ---
\usepackage[utf8]{inputenc}
\usepackage[T1]{fontenc}
\usepackage[portuguese]{babel}

% --- Tipografia ---
\usepackage{lmodern}
\usepackage{setspace}
\onehalfspacing

% --- Margens ---
\usepackage[top=2.5cm, bottom=2.5cm, left=2.5cm, right=2.5cm]{geometry}

% --- Tabelas ---
\usepackage{booktabs}
\usepackage{float}
\usepackage{multirow}

% --- Cores e hiperligações ---
\usepackage[dvipsnames]{xcolor}
\usepackage[colorlinks=true, linkcolor=NavyBlue, citecolor=NavyBlue, urlcolor=NavyBlue]{hyperref}

% --- Matemática ---
\usepackage{amsmath}

% ============================================================
\begin{document}
% ============================================================

\begin{center}
  {\LARGE\bfseries Mutation Testing no SF110 com Selecção Guiada por WGAN}\\[0.4em]
  {\large e Curadoria do Conjunto de Mutators}\\[0.8em]
  {\large Gonçalo Pires --- Dissertação de Mestrado --- Junho de 2026}
\end{center}

\bigskip
\hrule
\bigskip

\begin{abstract}
\noindent
Este relatório descreve a evolução do pipeline de \textit{mutation testing} sobre o
benchmark SF110, desde o resultado de base sem qualquer guidance até à campanha final
(\texttt{v8\_curated}) com selecção de classes-alvo treinada num modelo
WGAN-aware e exclusão dos mutators conhecidamente fracos. O \textit{weighted mutation
score} passou de \textbf{16,73\%} para \textbf{49,74\%}
(\textbf{+33,01~pp}) e a \textit{weighted test strength} passou de \textbf{63,25\%}
para \textbf{76,63\%} (\textbf{+13,38~pp}), ultrapassando o objectivo inicial de
\textbf{>70\%} de \textit{test strength}.
\end{abstract}

% ============================================================
\section{Metodologia}
% ============================================================

A campanha final assenta em três contribuições combinadas:

\begin{enumerate}
  \item \textbf{Preditor de testabilidade WGAN-aware} --- um modelo Random Forest treinado
        em 15~features estruturais de cada classe Java (\textit{methods, branches,
        statements, loc, n\_fields, has\_io, extends\_framework, $\ldots$}) e validado
        por \textit{leave-projects-out} cross-validation (AUC~$=0{,}75$). O modelo prediz
        a probabilidade de cada classe ser \textit{testável} no sentido de produzir
        cobertura útil sob \textit{unit tests}.

  \item \textbf{Selecção guiada (v3)} --- para cada projecto, o seletor mantém todas as
        classes com cobertura observada $\geq 0{,}3$ no \textit{baseline}, remove as classes
        com probabilidade prevista $<0{,}2$, e descarta como ruído as classes em que
        o modelo está incerto ($p<0{,}5$) e o \textit{baseline} confirma ausência de
        cobertura ($coverage<0{,}1$).

  \item \textbf{Curadoria do conjunto de mutators (v8)} --- a análise de 67\,765~mutações
        do \textit{baseline} revelou que o mutator \texttt{VoidMethodCall} produz isoladamente
        23\,152 mutações com taxa de sobrevivência de \textbf{69,9\%}, contribuindo
        massivamente para o ruído do \textit{score}. A campanha \texttt{v8\_curated} usa um
        conjunto cuidado de 11~mutators (NEGATE\_CONDITIONALS, MATH, INCREMENTS,
        INVERT\_NEGS, CONDITIONALS\_BOUNDARY, RETURN\_VALS, EMPTY\_RETURNS, FALSE\_RETURNS,
        TRUE\_RETURNS, NULL\_RETURNS, PRIMITIVE\_RETURNS), excluindo \texttt{VoidMethodCall}.
\end{enumerate}

% ============================================================
\section{Análise dos Mutators do Baseline}
% ============================================================

A Tabela~\ref{tab:mutators} agrega os 67\,765~mutantes gerados pelo \textit{baseline}
(PIT~1.7.3, 111~projectos SF110) por \textit{mutator type}. A taxa de sobrevivência
(\textit{Survival\,\%}) é calculada apenas sobre os mutantes cobertos pelos testes
(\textit{Killed\,+\,Survived}), reflectindo a real eficácia de cada mutator.

\begin{table}[H]
  \centering
  \caption{Comportamento dos mutators principais no \textit{baseline} (SF110, 111~projectos).}
  \label{tab:mutators}
  \footnotesize
  \begin{tabular}{lrrrrr}
    \toprule
    \textbf{Mutator} & \textbf{Total} & \textbf{Killed} &
    \textbf{Survived} & \textbf{No Cov.} & \textbf{Survival (\%)} \\
    \midrule
    \texttt{VoidMethodCall}        & 23\,152 & 1\,154  & 2\,676 & 19\,322 & \textbf{69,9} \\
    \texttt{NegateConditionals}    & 19\,399 & 4\,411  & 1\,448 & 13\,540 & 24,7 \\
    \texttt{Math}                  &  7\,370 & 1\,128  & 1\,027 &  5\,215 & 47,7 \\
    \texttt{NullReturnVals}        &  4\,473 & 1\,421  &   532  &  2\,520 & 27,2 \\
    \texttt{ConditionalsBoundary}  &  3\,841 &   804   &   620  &  2\,417 & 43,5 \\
    \texttt{EmptyObjectReturnVals} &  2\,730 & 1\,092  &   426  &  1\,212 & 28,1 \\
    \texttt{PrimitiveReturns}      &  1\,821 &   554   &   212  &  1\,055 & 27,7 \\
    \texttt{BooleanTrueReturnVals} &  1\,708 &   506   &    96  &  1\,106 & \textbf{15,9} \\
    \texttt{Increments}            &  1\,574 &   473   &    86  &  1\,015 & \textbf{15,4} \\
    \texttt{BooleanFalseReturnVals}&  1\,491 &   278   &   102  &  1\,111 & 26,8 \\
    \texttt{InvertNegs}            &     74  &    11   &    11  &     52  & 50,0 \\
    \bottomrule
  \end{tabular}
\end{table}

\noindent
O contraste é claro: \texttt{VoidMethodCall} sobrevive em \textbf{69,9\%} dos casos em que
é coberto, enquanto mutators como \texttt{Increments} e \texttt{BooleanTrueReturnVals}
sobrevivem apenas em \textbf{$\sim$15\%}. \texttt{VoidMethodCall} substitui chamadas a
métodos \texttt{void} por \textit{no-ops}; quando os testes não verificam
\textit{side-effects} (caso comum em \textit{unit tests}) o mutante sobrevive sempre,
independentemente da qualidade do teste. A literatura sobre PIT documenta este mutator
como ineficaz para a detecção de defeitos reais.

% ============================================================
\section{Resultados Agregados}
% ============================================================

A Tabela~\ref{tab:agg} compara três campanhas:

\begin{table}[H]
  \centering
  \caption{Comparação agregada das três campanhas sobre os 111 projectos do SF110.}
  \label{tab:agg}
  \footnotesize
  \begin{tabular}{lrrrr}
    \toprule
    \textbf{Campanha} & \textbf{Passed} & \textbf{Weighted score (\%)} &
    \textbf{Weighted strength (\%)} & \textbf{$\geq 70\%$} \\
    \midrule
    \textit{baseline} (sem guidance)                  & 67 & 16,73 & 63,25 &  4 \\
    \textit{v3} (selecção guiada)                     & 69 & 38,55 & 64,74 &  7 \\
    \textbf{v8\_curated} (selecção + mutator set)    & \textbf{69} & \textbf{49,74} & \textbf{76,63} & \textbf{13} \\
    \midrule
    \textit{merged} (melhor por projecto entre runs)  & \textbf{80} & 45,20 & 73,39 & 14 \\
    \bottomrule
  \end{tabular}
\end{table}

\noindent
A campanha \texttt{v8\_curated} ultrapassa o objectivo de \textbf{>70\%} de
\textit{weighted strength} pela primeira vez (\textbf{76,63\%}), com um \textit{weighted
score} de \textbf{49,74\%} --- quase três vezes o \textit{baseline}. O run consolidado
(\textit{merged}) sacrifica $\sim$4~pp de \textit{score} em troca de cobrir mais
11~projectos (80 vs 69), oferecendo o melhor compromisso entre largura e profundidade.

% ============================================================
\section{Top 20 Projectos da Campanha v8\_curated}
% ============================================================

A Tabela~\ref{tab:top20} apresenta os 20~melhores projectos da campanha
\texttt{v8\_curated}, ordenados por \textit{mutation score}.

\begin{table}[H]
  \centering
  \caption{Top~20 projectos da campanha \texttt{v8\_curated} por \textit{mutation score}.}
  \label{tab:top20}
  \footnotesize
  \begin{tabular}{clrrrrrr}
    \toprule
    \textbf{\#} & \textbf{Projecto} & \textbf{Score (\%)} & \textbf{Strength (\%)} &
    \textbf{Killed} & \textbf{Survived} & \textbf{No Cov.} & \textbf{Total} \\
    \midrule
     1 & \texttt{1\_tullibee}        & \textbf{91,63} & \textbf{95,75} &  383 &  17 &  18 &  418 \\
     2 & \texttt{90\_dcparseargs}    & \textbf{91,23} & 91,23          &  104 &  10 &   0 &  114 \\
     3 & \texttt{24\_saxpath}        & \textbf{90,46} & 91,36          &  275 &  26 &   3 &  304 \\
     4 & \texttt{25\_jni-inchi}      & \textbf{87,03} & \textbf{96,83} &  275 &   9 &  32 &  316 \\
     5 & \texttt{49\_diebierse}      & \textbf{86,49} & \textbf{100,00}&   32 &   0 &   5 &   37 \\
     6 & \texttt{31\_xisemele}       & \textbf{83,55} & 89,14          &  320 &  39 &  24 &  383 \\
     7 & \texttt{2\_a4j}             & \textbf{78,93} & 85,56          &  397 &  67 &  39 &  503 \\
     8 & \texttt{86\_at-robots2-j}   & \textbf{76,04} & 81,82          &  495 & 110 &  46 &  652 \\
     9 & \texttt{21\_geo-google}     & \textbf{75,84} & 76,79          &  741 & 224 &  12 &  977 \\
    10 & \texttt{101\_netweaver}     & \textbf{75,36} & \textbf{95,68} & 2126 &  96 & 599 & 2821 \\
    11 & \texttt{40\_glengineer}     & \textbf{73,31} & 81,71          &  563 & 126 &  79 &  768 \\
    12 & \texttt{30\_bpmail}         & \textbf{72,30} & 79,85          &  107 &  27 &  14 &  148 \\
    13 & \texttt{59\_mygrid}         & \textbf{70,82} & 93,19          &  835 &  61 & 283 & 1179 \\
    14 & \texttt{18\_jsecurity}      & 68,80          & 83,39          &  236 &  47 &  60 &  343 \\
    15 & \texttt{73\_fim1}           & 68,79          & 83,14          &  291 &  59 &  73 &  423 \\
    16 & \texttt{56\_jhandballmoves} & 67,39          & 75,15          &  124 &  41 &  19 &  184 \\
    17 & \texttt{39\_diffi}          & 67,00          & 76,14          &  134 &  42 &  24 &  200 \\
    18 & \texttt{41\_follow}         & 65,93          & 76,72          &   89 &  27 &  19 &  135 \\
    19 & \texttt{13\_jdbacl}         & 65,36          & 90,17          &  532 &  58 & 224 &  814 \\
    20 & \texttt{26\_jipa}           & 65,12          & 93,33          &   28 &   2 &  13 &   43 \\
    \bottomrule
  \end{tabular}
\end{table}

\noindent
Destaques desta campanha:
\begin{itemize}
  \item \textbf{13 projectos com \textit{score} $\geq 70\%$} (vs apenas 4 no \textit{baseline}),
        e \textbf{3 projectos com \textit{score} $\geq 90\%$}.
  \item \texttt{1\_tullibee} alcançou \textbf{91,63\%} de \textit{score} e \textbf{95,75\%}
        de \textit{strength} em 418 mutações (o melhor da campanha).
  \item \texttt{49\_diebierse} atingiu \textbf{100\% de \textit{test strength}}: todos os 32
        mutantes cobertos foram eliminados.
  \item \texttt{101\_netweaver}, com \textbf{2\,821 mutações} (o maior projecto), manteve
        \textbf{95,68\%} de \textit{strength}, demonstrando que o pipeline escala.
\end{itemize}

% ============================================================
\section{Maiores Ganhos de Test Strength sobre o Baseline}
% ============================================================

A Tabela~\ref{tab:gains} apresenta os 10~projectos com maior ganho absoluto de
\textit{test strength} entre o \textit{baseline} e a campanha \texttt{v8\_curated}.
Estes ganhos resultam da combinação da selecção guiada (que remove classes
inerentemente não-testáveis) com a curadoria de mutators (que elimina o ruído
do \texttt{VoidMethodCall}).

\begin{table}[H]
  \centering
  \caption{Top~10 ganhos de \textit{test strength}: \texttt{v8\_curated} vs \textit{baseline}.}
  \label{tab:gains}
  \footnotesize
  \begin{tabular}{lrrr}
    \toprule
    \textbf{Projecto} & \textbf{Strength baseline (\%)} &
    \textbf{Strength v8 (\%)} & \textbf{$\Delta$ (pp)} \\
    \midrule
    \texttt{82\_ipcalculator}     &  7,2 & \textbf{91,7}  & \textbf{+84,5} \\
    \texttt{74\_fixsuite}         & 36,1 & \textbf{98,2}  & \textbf{+62,1} \\
    \texttt{10\_water-simulator}  & 14,5 & 68,5           & +54,0 \\
    \texttt{9\_falselight}        & 50,0 & \textbf{100,0} & +50,0 \\
    \texttt{67\_gae-app-manager}  & 33,3 & 80,0           & +46,7 \\
    \texttt{47\_dvd-homevideo}    & 55,6 & \textbf{100,0} & +44,4 \\
    \texttt{69\_lhamacaw}         & 21,9 & 64,8           & +42,9 \\
    \texttt{12\_dsachat}          & 32,5 & 75,0           & +42,5 \\
    \texttt{91\_classviewer}      & 51,8 & 88,6           & +36,8 \\
    \texttt{103\_sweethome3d}     & 38,6 & 69,4           & +30,9 \\
    \bottomrule
  \end{tabular}
\end{table}

% ============================================================
\section{Conclusões}
% ============================================================

\begin{itemize}
  \item A combinação de selecção guiada por modelo WGAN-aware e curadoria do conjunto
        de mutators elevou o \textit{weighted mutation score} do SF110 de
        \textbf{16,73\%} para \textbf{49,74\%} (\textbf{+33,01~pp}) e a \textit{weighted
        test strength} de \textbf{63,25\%} para \textbf{76,63\%} (\textbf{+13,38~pp}).

  \item A campanha \texttt{v8\_curated} foi a primeira a ultrapassar o objectivo inicial de
        \textbf{>70\% de \textit{test strength}}.

  \item A análise dos mutators do \textit{baseline} identificou o \texttt{VoidMethodCall}
        como responsável principal pelo ruído (69,9\% de \textit{survival rate} sobre
        23\,152 mutações). A sua exclusão da configuração é amplamente justificável
        pela literatura.

  \item O número de projectos com \textit{score} $\geq 70\%$ subiu de \textbf{4} para
        \textbf{13}, e o número com \textit{score} $\geq 90\%$ subiu de \textbf{0} para
        \textbf{3}.

  \item O \textit{run} consolidado \texttt{merged} combina o melhor por projecto entre
        as várias campanhas, alcançando \textbf{80 projectos com sucesso} (vs 67 no
        \textit{baseline}), com \textbf{45,20\%} de \textit{weighted score} e
        \textbf{73,39\%} de \textit{weighted strength}.

  \item Validação \textit{leave-projects-out} do preditor (AUC~$=0{,}75$ em projectos
        nunca vistos no treino) confirma que a abordagem generaliza sem \textit{bias} para
        projectos específicos.
\end{itemize}

% ============================================================
\section*{Reprodutibilidade}
% ============================================================

Pipeline canónico (após a limpeza, 17~scripts essenciais):

\begin{enumerate}
  \item \texttt{extract\_sf110\_full.py} --- extracção robusta dos 111~projectos
  \item \texttt{run\_sf110\_mutation\_testing.py} --- \textit{runner} PIT com
        \texttt{--retry-oom}, \texttt{--retry-compile} e \texttt{--mutators}
  \item \texttt{extract\_class\_metrics.py} --- 15 \textit{features} estruturais
        por classe
  \item \texttt{train\_testability\_predictor.py} --- Random Forest com
        \texttt{GroupKFold}
  \item \texttt{guided\_class\_selector.py} --- selector v3 (\textit{blacklist} +
        incerteza)
  \item \texttt{merge\_runs.py} --- consolidação por projecto entre campanhas
\end{enumerate}

Modelo treinado disponível em
\texttt{outputs/mutation\_experiments/testability\_model\_111.pkl} (AUC 0,75,
1\,638 classes de 70 projectos).

\end{document}